\cvheader{\CandidateCityEN}{Agricultural Engineer | Dairy Scientist | Applied Livestock Research}

\section*{Profile}
% CUSTOMIZE: Mirror the most important requirements from the job description.
Agricultural engineer and dairy scientist with 4+ years of experience in applied livestock research, experimental design, project coordination and statistical analysis. Doctoral research develops scalable, low-cost gas-sensing systems for livestock emissions monitoring. Experienced in commercial dairy-farm studies, integration of herd-management and environmental data, and communication of scientific results. Seeking to translate rigorous research into practical innovations for ruminant production.

\section*{Professional Experience}
\entry{Doctoral Researcher (PhD Candidate)}{Mar 2023--present}{Leibniz Institute for Agricultural Engineering and Bioeconomy (ATB), Potsdam}{Freie Universität Berlin}
\begin{cvitems}
  \item Plan and conduct interdisciplinary field research in commercial dairy housing; coordinate experimental work with scientists, technicians and interns.
  \item Develop and validate scalable low-cost gas-sensing systems using NDIR and electrochemical sensors against reference-grade analytical methods.
  \item Analyse long-term data in R and integrate emissions and barn-climate measurements with herd records, including milk yield, reproduction and health data.
  \item Prepare publications, technical documentation and conference presentations; supervise interns and coordinate field campaigns.
\end{cvitems}

\entry{Scientist (full-time)}{Feb 2022--Feb 2023}{Leibniz Institute for Agricultural Engineering and Bioeconomy (ATB), Potsdam}{}
\begin{cvitems}
  \item Planned and executed applied research on barn climate and gaseous emissions under practical dairy-farm conditions.
  \item Installed, operated and validated FTIR, CRDS, GC-MS, TDLAS, PAS and environmental monitoring systems; evaluated complex datasets and documented results.
\end{cvitems}

\entry{Research Assistant / Master's Thesis (part-time)}{Apr 2021--Jan 2022}{Leibniz Institute for Agricultural Engineering and Bioeconomy (ATB), Potsdam}{}
\begin{cvitems}
  \item Investigated the spatial distribution of methane, ammonia and carbon dioxide in a naturally ventilated dairy barn, from experimental setup through statistical evaluation.
\end{cvitems}

\section*{Education}
\entry{PhD in Agricultural Engineering (ongoing)}{}{Freie Universität Berlin}{}
\entry{M.Sc. Dairy Science, grade 1.6}{}{Christian-Albrechts-Universität zu Kiel}{}
Relevant study areas: ruminant nutrition, herd health, precision livestock farming, forage quality, eco-efficiency, dairy economics and biometrical planning.

\entry{B.Tech. Dairy Technology}{2015--2019}{College of Dairy Science and Food Technology, India}{}

\section*{Evidence \& Expertise}
\textbf{Research output:} 1 first-author peer-reviewed article; 1 manuscript under review; 8 presentations across 9 scientific conferences; work on 4 research projects.\par
\textbf{Methods:} Study planning; statistical analysis (R); scientific writing; literature review; project coordination; R, Python, MATLAB and LaTeX.\par
\textbf{Dairy and R\&D:} Ruminant nutrition; herd data; livestock housing; environmental monitoring; sensor development and field validation.

\section*{Leadership, Languages \& Practical Information}
\textbf{Leadership:} Co-Chair, EurAgEng Young Professionals Network; PhD Representative, ATB; Communication Manager, Leibniz PhD Network.\par
\textbf{Languages:} Hindi (native); German (B2); English (B2). \quad
\textbf{Work authorisation:} Permanent residence in Germany; no sponsorship required.\par
\textbf{Driving licence:} Indian Category B and International Driving Permit; German Class B practical examination in progress.\par
\textbf{References:} Prof. Dr. Thomas Amon and Dr.-Ing. David Janke, ATB.

